Forecasting rising memory pressure can support monitoring, but accurate forecasts do not by themselves identify memory leaks. We compare persistence, linear trend, restricted ARIMA and a small LSTM on one Antarex fault-injection recording and three Westermo industrial recordings, using chronological splits and 60- or 300-second horizons. Forecast error and growth-warning performance are evaluated separately. In the original Antarex run, LSTM reduces RMSE from ARIMA's 1.516 to 1.216 percentage points of RAM. Its warning rule, however, recognizes 34 of 35 injected events against trend-only's 35, with timestamp recall 0.466 versus 0.859. A matched study of five random seeds and two LSTM input configurations adds 40 runs. Extra features improve mean RMSE on only two of four recordings. Multivariate Antarex warning F1 is $0.571\pm0.031$, below trend-only's 0.830, and false alarm starts range from 43 to 137 against 86. Analyzing the saved scores shows that the original missed event occurred simply because the forecast-growth threshold was not reached. Because the Westermo dataset lacks verified leak ground truth, it can only support forecasting evaluations. Furthermore, while local replay and inference benchmarks confirm prototype viability, we have not yet validated live runtime performance or lead-time warnings prior to memory exhaustion.
